{\rtf1\ansi\ansicpg1252\cocoartf2639
\cocoatextscaling0\cocoaplatform0{\fonttbl\f0\fswiss\fcharset0 Helvetica;\f1\fnil\fcharset0 LucidaGrande;\f2\fnil\fcharset0 Menlo-Regular;
}
{\colortbl;\red255\green255\blue255;}
{\*\expandedcolortbl;;}
\paperw11900\paperh16840\margl1440\margr1440\vieww11520\viewh8400\viewkind0
\pard\tx566\tx1133\tx1700\tx2267\tx2834\tx3401\tx3968\tx4535\tx5102\tx5669\tx6236\tx6803\pardirnatural\partightenfactor0

\f0\fs24 \cf0 \\documentclass[11pt]\{article\}\
\
% ===================== PACKAGES =====================\
\\usepackage[a4paper,margin=1in]\{geometry\}\
\\usepackage\{amsmath,amssymb\}\
\\usepackage\{enumitem\}\
\\usepackage\{fancyhdr\}\
\\usepackage\{xcolor\}\
\
% ===================== HEADER & FOOTER =====================\
\\pagestyle\{fancy\}\
\\fancyhf\{\}\
\\lhead\{CSE 400: Fundamentals of Probability in Computing\}\
\\rhead\{Lecture Scribe -- Lecture 10\}\
\\cfoot\{\\thepage\}\
\
% ===================== TITLE =====================\
\\title\{\
    \\normalsize School of Engineering and Applied Science (SEAS), Ahmedabad University \\\\\
    \\vspace\{0.2cm\}\
    \\textbf\{CSE 400: Fundamentals of Probability in Computing\}\\\\\
    \\Large Lecture 10 Scribe: Randomized Min-Cut Algorithm\
\}\
\
\\author\{\}\
\\date\{\}    \
\\begin\{document\}\
\\maketitle\
\
\
\\vspace\{-2cm\}\
\\begin\{center\}\
    \\begin\{tabular\}\{ll\}\
        \\textbf\{Instructor:\} & \{Prof. Dhaval Patel\} \\\\ [1.5ex]\
        \\textbf\{Course:\} & \{CSE 400 -- Fundamentals of Probability in Computing\} \\\\ [1.5ex]\
        \\textbf\{Lecture:\} & \{Lecture 10\} \\\\ [1.5ex]\
        \\textbf\{Date:\} & \{5th Feb 2026\} \\\\ [1.5ex]\
    \\end\{tabular\}\
\\end\{center\}\
\
\\hrule\
\\vspace\{0.5cm\}\
\
% ===================== LECTURE SCRIBE =====================\
\
\\section\{Overview of the Lecture\}\
This lecture covers:\
\\begin\{itemize\}\
    \\item Min-Cut Problem\
    \\item Motivation for using min-cut\
    \\item Definition of min-cut\
    \\item Successful and unsuccessful min-cut runs\
    \\item Max-Flow Min-Cut Theorem\
    \\item Deterministic Min-Cut Algorithm\
    \\item Stoer--Wagner minimum cut algorithm\
    \\item Randomized Min-Cut Algorithm\
    \\item Karger\'92s Randomized Algorithm\
    \\item Comparison between deterministic and randomized approaches\
    \\item Theorem for min-cut set\
    \\item Python simulation\
\\end\{itemize\}\
\
\\section\{Min-Cut Problem\}\
\
\\subsection\{Why Use Min-Cut?\}\
Min-cut algorithms are used to solve problems related to:\
\\begin\{itemize\}\
    \\item Network connectivity\
    \\item Reliability\
    \\item Optimization\
\\end\{itemize\}\
\
Applications include:\
\\begin\{itemize\}\
    \\item \\textbf\{Network Design:\} Improves communication efficiency and helps find the minimum capacity cut.\
    \\item \\textbf\{Communication Networks:\} Helps understand vulnerability to failures and build fault-tolerant networks.\
    \\item \\textbf\{VLSI Design:\} Used for partitioning circuits into smaller components to reduce interconnectivity complexity.\
\\end\{itemize\}\
\
\\subsection\{Definition of Min-Cut\}\
\
\\subsubsection\{Cut-Set\}\
A cut-set in a graph is a set of edges whose removal breaks the graph into two or more connected components.\
\
\\subsubsection\{Minimum Cut\}\
Given a graph \\( G = (V, E) \\) with \\( n \\) vertices, the minimum cut (min-cut) problem is to find a minimum cardinality cut-set in \\( G \\).\
\
Min-cut algorithms such as Karger\'92s algorithm are random and can be sensitive to the initial choice of edges. Contracting critical edges early may result in an incorrect cut.\
\
\\section\{Edge Contraction\}\
Edge contraction is the primary operation used in min-cut algorithms.\
\
\\textbf\{Definition:\}  \
Edge contraction removes an edge from a graph while simultaneously merging the two vertices.\
\
\\textbf\{Procedure:\}  \
For an edge \\( (u,v) \\):\
\\begin\{itemize\}\
    \\item Merge vertices \\( u \\) and \\( v \\)\
    \\item Remove all edges between \\( u \\) and \\( v \\)\
    \\item Retain all other edges\
\\end\{itemize\}\
\
The resulting graph may contain parallel edges but no self-loops.\
\
\\section\{Successful and Unsuccessful Min-Cut Runs\}\
\
\\subsection\{Successful Min-Cut Run\}\
A successful min-cut run refers to an execution of a min-cut algorithm that correctly identifies the minimum cut of the graph.\
\
\\subsection\{Unsuccessful Min-Cut Run\}\
An unsuccessful min-cut run refers to an execution in which the algorithm incorrectly identifies the minimum cut due to contraction of critical edges.\
\
\\section\{Max-Flow Min-Cut Theorem\}\
\
\\textbf\{Statement:\}  \
In a flow network, the maximum amount of flow from the source to the sink is equal to the total weight of the edges in a minimum cut.\
\
\\textbf\{Definitions:\}\
\\begin\{itemize\}\
    \\item \\textbf\{Capacity of a cut:\} Sum of capacities of edges oriented from \\( X \\) to \\( Y \\)\
    \\item \\textbf\{Minimum cut:\} Cut with the smallest capacity\
    \\item \\textbf\{Minimum cut capacity:\} Capacity of the minimum cut\
    \\item \\textbf\{Maximum flow:\} Largest possible flow from source \\( S \\) to sink \\( T \\)\
\\end\{itemize\}\
\
\\section\{Deterministic Min-Cut Algorithm\}\
\
\\subsection\{Stoer--Wagner Minimum Cut Algorithm\}\
Let \\( s \\) and \\( t \\) be two vertices of a graph \\( G \\). Let \\( G/\\\{s,t\\\} \\) be the graph obtained by merging \\( s \\) and \\( t \\).\
\
A minimum cut of \\( G \\) is the smaller of:\
\\begin\{itemize\}\
    \\item A minimum \\( s\\!-\\!t \\) cut of \\( G \\)\
    \\item A minimum cut of \\( G/\\\{s,t\\\} \\)\
\\end\{itemize\}\
\
\\subsection\{Stoer--Wagner Algorithm Pseudocode\}\
\
\\begin\{verbatim\}\
Algorithm 1: MinimumCutPhase(G, a)\
A 
\f1 \uc0\u8592 
\f0  \{a\};\
while A \uc0\u8800  V do\
    add to A the most tightly connected vertex;\
return the cut weight as the cut of the phase;\
\
Algorithm 2: MinimumCut(G)\
while |V| \uc0\u8805  1 do\
    choose any a from V;\
    MinimumCutPhase(G, a);\
    if cut-of-the-phase is lighter than current minimum cut then\
        store it;\
    shrink G by merging the last two added vertices;\
return the minimum cut;\
\\end\{verbatim\}\
\
\\section\{Randomized Min-Cut Algorithm\}\
\
\\subsection\{Why Randomized Algorithms?\}\
Randomized algorithms provide a probabilistic guarantee of success and often require fewer iterations. They offer:\
\\begin\{itemize\}\
    \\item Efficiency\
    \\item Parallelization\
    \\item Approximation guarantees\
    \\item Robustness\
\\end\{itemize\}\
\
\\subsection\{Karger\'92s Randomized Algorithm\}\
Karger\'92s algorithm repeatedly contracts randomly selected edges until two vertices remain.\
\
\\subsection\{Recursive Randomized Min-Cut Pseudocode\}\
\
\\begin\{verbatim\}\
else\
for i 
\f1 \uc0\u8592 
\f0  1 to \uc0\u945  do\
    G' 
\f1 \uc0\u8592 
\f0  multigraph obtained by applying\
         n \uc0\u8722  
\f2 \uc0\u8968 
\f0 n/\uc0\u8730 \u945 
\f2 \uc0\u8969 
\f0  random contraction steps;\
    C' 
\f1 \uc0\u8592 
\f0  Recursive-Randomized-Min-Cut(G', \uc0\u945 );\
    if i = 1 or |C'| < |C| then\
        C 
\f1 \uc0\u8592 
\f0  C';\
return C;\
\\end\{verbatim\}\
\
\\section\{Comparison: Deterministic vs Randomized Min-Cut\}\
\
\\subsection\{Deterministic Min-Cut\}\
\\begin\{itemize\}\
    \\item Guarantees exact minimum cut\
    \\item Higher time complexity for large graphs\
    \\item Time complexity:\
    \\[\
    O(V \\cdot E + V^2 \\log V)\
    \\]\
\\end\{itemize\}\
\
\\subsection\{Randomized Min-Cut\}\
\\begin\{itemize\}\
    \\item Produces minimum cut with high probability\
    \\item Time complexity:\
    \\[\
    O(V^2)\
    \\]\
\\end\{itemize\}\
\
\\section\{Theorem for Min-Cut Set\}\
The algorithm outputs a min-cut set with probability at least:\
\\[\
\\frac\{2\}\{n(n-1)\}\
\\]\
\
\\section\{Python Simulation\}\
\\begin\{itemize\}\
    \\item Open Campuswire post for Lecture 10\
    \\item Download the provided \\texttt\{.ipynb\} file\
\\end\{itemize\}\
\
\\section\{Summary\}\
Key concepts for revision:\
\\begin\{itemize\}\
    \\item Cut-set and min-cut definitions\
    \\item Edge contraction\
    \\item Successful vs unsuccessful runs\
    \\item Max-Flow Min-Cut theorem\
    \\item Stoer--Wagner algorithm\
    \\item Karger\'92s randomized algorithm\
    \\item Probability guarantee for min-cut\
\\end\{itemize\}\
\
\\bigskip\
\\hrule\
\\vspace\{0.2cm\}\
\\begin\{center\}\
    \\small \\textit\{End of Lecture Scribe\}\
\\end\{center\}\
\
\\end\{document\}\
}